%----------
%	CONTENIDO NEGATIVO - LLMs
%----------	


In this Section, we applied Generative AI to the "Contenido Negativo" task, conducting a total of 64 experiments across 8 different windows, as detailed in Table~\ref{tab:window-metrics-cnegativo-genai}. As explained in the corresponding section for ML (Section~\ref{ageneral-genai}), W1-W4 represented experiments utilizing "Context 0", while W5-W8 corresponded to those using "Context 1". Within each context, the first window (W1 and W5) used "Prompt\_0\_EN", the second (W2 and W6) employed "Prompt\_0\_ES", the third (W3 and W7) utilized "Prompt\_1\_EN", and the final window (W4 and W8) applied "Prompt\_1\_ES".


\begin{table}[H]
    \centering
    \begingroup
    \small % Change to desired font size
    \ttabbox[\FBwidth]{
        \caption{CNEG Average Evaluation Metrics per Window in GenAI Experiments}
        \label{tab:window-metrics-cnegativo-genai} % Unique label for the table
    }{
        \begin{tabularx}{\textwidth}{
            X
            X
            X
            >{\centering\arraybackslash}X
            >{\centering\arraybackslash}X
            >{\columncolor{gray!15}\centering\arraybackslash}X % Highlight this column
            >{\centering\arraybackslash}X
            >{\centering\arraybackslash}X
            >{\centering\arraybackslash}X
            >{\centering\arraybackslash}X
        }
            \hline
            \rowcolor{gray!30} % Adding a gray color to the header row
            \textbf{} & \textbf{Context} & \textbf{Prompt} & \textbf{Acc. (Std.)} & \textbf{Time (s)} & \textbf{W. Avg. F1-Score} & \textbf{D. Víc. F1-Score} & \textbf{D. Acto F1-Score} & \textbf{Ins. F1-Score} & \textbf{D. D. Aut. F1-Score} \\
            \hline
            W1 & 0 & 0\_EN & 0.248 (0.136) & 138.517 & 0.317 & 0.329 & 0.180 & 0.336 & 0.042 \\
            \hline
            W2 & 0 & 0\_ES & 0.247 (0.101) & 166.576 & 0.296 & 0.338 & 0.212 & 0.295 & 0.085 \\
            \hline
            W3 & 0 & 1\_EN & 0.248 (0.139) & 156.534 & 0.322 & 0.338 & 0.203 & 0.324 & 0.021 \\
            \hline
            W4 & 0 & 1\_ES & 0.258 (0.121) & 154.562 & 0.296 & 0.330 & 0.224 & 0.304 & 0.131 \\
            \hline
            W5 & 1 & 0\_EN & 0.249 (0.096) & 159.723 & 0.321 & 0.349 & 0.145 & 0.328 & 0.051 \\
            \hline
            W6 & 1 & 0\_ES & 0.236 (0.181) & 153.319 & 0.312 & 0.338 & 0.074 & 0.275 & 0.071 \\
            \hline
            W7 & 1 & 1\_EN & 0.263 (0.121) & 157.551 & 0.330 & 0.339 & 0.166 & 0.360 & 0.077 \\
            \hline
            W8 & 1 & 1\_ES & 0.265 (0.127) & 161.590 & 0.327 & 0.339 & 0.184 & 0.359 & 0.084 \\
            \hline
            \multicolumn{10}{l}{Source: 'Training\_GenAI.xlsx'} \\
            \hline
        \end{tabularx}
    }
    \endgroup
\end{table}


Figure~\ref{fig:exp_contenido_negativo_genai} shows the distribution of experiments for "Contenido Negativo" using the Generative AI approach. Overall, we observed a positive trend across the experiments, with an $R^2$ value of 0.18, although for Context 0 an irregular pattern was seen where W1 (0.317) and W3 (0.322), corresponding to prompts in English, outperformed W2 (0.296) and W4 (0.296), corresponding to prompts in Spanish, in terms of weighted average F1-Score. However, for Context 1 we saw a more linear increasing pattern, except for W6, achieving the best overall weighted average F1-Score of 0.330 in W7, which corresponded to English prompt 1. Moreover, it was observed that for odd-numbered windows (W1, W3, W5, W7) the performance in terms of accuracy and weighted average F1-Score was slightly better, with a mean weighted average F1-Score of 0.319 and 0.304 for "Prompt\_0\_EN" and "Prompt\_0\_ES", respectively, and 0.326 for "Prompt\_1\_EN" and 0.311 for "Prompt\_1\_ES".

Figure~\ref{fig:exp_contenido_negativo_classes_genai} illustrates the distribution of experiments across different classes. Notably, there was a positive trend for the classes "Desprestigiar Víctima", "Insultos", and "Desprestigiar Deportista Autora", while "Desprestigiar Acto" showed a negative trend. When comparing contexts, Context 0 (W1-W4) and Context 1 (W5-W8) generally performed similarly, with the exception of "Desprestigiar Acto", where the mean weighted average F1-Score was significantly higher in Context 0 (0.205) than in Context 1 (0.142). A similar pattern was observed with different prompts, as the most notable difference occurred between "prompt\_0\_EN" and "prompt\_0\_ES" (0.046 vs. 0.078), and "prompt\_1\_EN" and "prompt\_1\_ES" (0.049 vs. 0.108) for the "Desprestigiar Deportista Autora" class. These results suggest that the varying contexts and prompts did not significantly influence the classification results.



\begin{figure}[H]
	\ffigbox[\FBwidth]
	{\caption{Distribution of the GenAI Experiments over Time for "Contenido Negativo" with the Different Windows and Weighted Avg. F1-Score trend line}\label{fig:exp_contenido_negativo_genai}}
 % [Aquí ponemos como queremos que aparezca en el índice, sin la referencia]{Aquí como queremos que aparezca en el documento, con la referencia}
	{\includegraphics[scale=0.25]{Plantilla_TFG_ingles_2019/imagenes/GenAI_Experiments in _ContenidoNegativo_ weighted.pdf}}
\end{figure}


\begin{figure}[H]
	\ffigbox[\FBwidth]
	{\caption{Distribution of the GenAI Experiments over Time for "Contenido Negativo", per Class with the Different Windows, and Trend Lines for Each Class}\label{fig:exp_contenido_negativo_classes_genai}}
 % [Aquí ponemos como queremos que aparezca en el índice, sin la referencia]{Aquí como queremos que aparezca en el documento, con la referencia}
	{\includegraphics[scale=0.25]{Plantilla_TFG_ingles_2019/imagenes/GenAI_Experiments in _ContenidoNegativo_f1score.pdf}}
\end{figure}


Table~\ref{tab:genai-cnegativo-genai-bestmodels_metrics} summarizes the top 5 best-performing generative models for the "Contenido Negativo" task. The weighted average F1-Scores ranged from 0.337 to 0.341, while accuracy was from 0.259 to 0.278. Overall, these models demonstrated better performance in the "Insultos" category and weaker performance in the "Desprestigiar Deportista Autora" category. A possible reason for this is that "Insultos" is a more general class, which was likely better understood by the pre-trained language models, whereas "Desprestigiar Víctima", "Desprestigiar Acto" and "Desprestigiar Deportista Autora" are more specialized categories, specifically created for this classification task.


\begin{table}[H]
    \centering
    \begingroup
    \small % Change to desired font size
    \ttabbox[\FBwidth]{
        \caption{CNEG Evaluation Metrics for the Best GenAI Models}
        \label{tab:genai-cnegativo-genai-bestmodels_metrics} % Unique label for the table
    }{
        \begin{tabularx}{\textwidth}{
            X
            >{\centering\arraybackslash}X
            >{\columncolor{gray!15}\centering\arraybackslash}X % Highlight this column
            >{\centering\arraybackslash}X
            >{\centering\arraybackslash}X
            >{\centering\arraybackslash}X
            >{\centering\arraybackslash}X
        }
            \hline
            \rowcolor{gray!30} % Adding a gray color to the header row
            \textbf{} & \textbf{Acc. (Std.)} & \textbf{Weighted Avg. F1-Score} & \textbf{D. Víc. F1-Score} & \textbf{D. Acto F1-Score} & \textbf{Ins. F1-Score} & \textbf{D. D. Aut. F1-Score} \\
            \hline
            \multicolumn{7}{l}{\textbf{Naive Models}} \\
            \hline
            \cnegmajor & 0.421 (0.027) & 0.250 & 0.000 & 0.000 & 0.000 & 0.592 \\
            \hline
            \cnegrandom & 0.289 (0.024) & 0.289 & 0.201 & 0.205 & 0.169 & 0.420 \\
            \hline
            \multicolumn{7}{l}{\textbf{Zero-Shot Models}} \\
            \hline
            \cneggemmatwob & 0.278 (0.146) & 0.341 & 0.347 & 0.186 & 0.382 & 0.099 \\
            \hline
            \cnegmistralsevenb & 0.259 (0.128) & 0.340 & 0.346 & 0.198 & 0.355 & 0.013 \\
            \hline
            \cnegllamathreeonea & 0.259 (0.158) & 0.340 & 0.352 & 0.205 & 0.342 & 0.013 \\
            \hline
            \cnegllamathreeonec & 0.273 (0.167) & 0.339 & 0.345 & 0.107 & 0.398 & 0.111  \\
            \hline
            \cnegllamathreeoneb & 0.267 (0.091) & 0.337 & 0.342 & 0.157 & 0.376 & 0.074 \\
            \hline
            \multicolumn{7}{l}{Source: 'Training\_GenAI.xlsx'} \\
            \hline
        \end{tabularx}
    }
    \endgroup
\end{table}

\newpage

Therefore, we selected model \textbf{\cneggemmatwob} as the best-performing model for this task and approach, due to achieving the highest weighted average F1-Score and accuracy. It is important to note that compared to the naive models, {\cnegmajor} only outperformed in accuracy, which in this case does not indicate a better model as the F1-Scores for the individual classes were 0, unless for the majority ("Desprestigiar Deportista Autora"). Moreover, {\cnegrandom} was better only in "Desprestigiar Acto" and "Desprestigiar Deportista Autora" F1-Scores, but the reason is that such model predicts at random and therefore approximating the probability of each class, but in the Zero-Shot learning approach was not relevant, as the dataset was not shown to the model. In Figure~\ref{fig:exp_cnegativo_genai_bestmodel} we present the confusion matrix for this model, where 68 instances of "Desprestigiar Víctima" and 63 instances of "Insultos" were incorrectly classified. Additionally, "Desprestigiar Acto" had some confusion with "Desprestigiar Víctima," indicating that the model might be confusing these two classes due to their conceptual similarities. Consequently, we can conclude that while \textbf{\cneggemmatwob} performed reasonably well in classes with general meaning ("Insultos"), there remains considerable room for improvement in distinguishing between the more specific classes within "Contenido Negativo".


% Best Model
\begin{figure}[H]
	\ffigbox[\FBwidth]
	{\caption{Confusion Matrix for the Best GenAI Model (\textbf{\cneggemmatwob}) for "Contenido Negativo"}\label{fig:exp_cnegativo_genai_bestmodel}}
 % [Aquí ponemos como queremos que aparezca en el índice, sin la referencia]{Aquí como queremos que aparezca en el documento, con la referencia}
	{\includegraphics[scale=0.11]{Plantilla_TFG_ingles_2019/imagenes/GenAI_Experiments in _ContenidoNegativo_ best_model.pdf}}
\end{figure}

\newpage